\section{Final individual project}

The final individual project accounts for 65\% of the overall course grade. Students design and execute an independent Computational Communication Science project that demonstrates their ability to formulate a research question, collect and analyse data, and apply fundamental techniques taught in the course.

This is an individual assessment. The report and all submitted code must be the student's own work, except where external sources or permitted generative-AI contributions are appropriately disclosed in accordance with the course manual.

Your submission must include:
\begin{enumerate}
    \item a PDF report; and
    \item a \texttt{.zip} folder containing all relevant code and output files.
\end{enumerate}

\subsection*{PDF report}

The recommended length of the report is five to six pages. It should include the following components.

\begin{enumerate}[A.]
    \item \textbf{Introduction} (one to two pages)
    \begin{itemize}
        \item A clearly formulated research question;
        \item Any subquestions and/or hypotheses;
        \item References to relevant literature, including course readings where appropriate.
    \end{itemize}

    \item \textbf{Analytic strategy}
    \begin{itemize}
        \item A transparent overview of the methodological approach;
        \item A justification for the choice of techniques, referring explicitly to methods discussed during the course.
    \end{itemize}

    \item \textbf{Dataset description}
    \begin{itemize}
        \item A carefully collected and relevant dataset of non-trivial size, for example retrieved via APIs or constructed by combining existing datasets;
        \item An exploratory description of the dataset;
        \item A brief argument explaining why the dataset is suitable for the research question.
    \end{itemize}

    \item \textbf{Conclusion and critical reflection}
    \begin{itemize}
        \item A well-substantiated answer to the research question;
        \item A critical discussion of the strengths, limitations, and uncertainties of the research design, data, analysis, and conclusions;
        \item An explanation of the implications of these limitations for the interpretation of the findings, as well as feasible directions for future research.
    \end{itemize}
\end{enumerate}

The conclusion and critical-reflection component contributes to the assessment of ILO H: students justify methodological choices, acknowledge limitations and uncertainties, and respond constructively to critical questions about their own research.

\subsection*{Project files (\texttt{.zip} folder)}

Alongside the PDF report, include a compressed folder containing:
\begin{itemize}
    \item Python scripts or notebooks for data collection, preprocessing, analysis, and visualisation. These files must be well documented, reproducible, and tailored to the student's own project;
    \item Output files, such as tables, plots, or summary statistics;
    \item A \texttt{README} file describing:
    \begin{itemize}
        \item how to run the code;
        \item what each script or notebook does; and
        \item any software or library requirements.
    \end{itemize}
\end{itemize}

\subsection*{Requirements and guidance}

Depending on the chosen topic, students must apply at least one technique taught in the course. Projects should include an exploratory description of the dataset and apply fundamental analytical techniques.

During lab sessions and individual consultations, students may discuss:
\begin{itemize}
    \item the scope of the project;
    \item the methodological requirements that the project must meet; and
    \item how the chosen analytical methods contribute to the final grade.
\end{itemize}

\textbf{Important notes}
\begin{itemize}
    \item All external data sources and code, including libraries, templates, and tutorials, must be properly cited.
    \item Code must be tailored to the student's own project and run without errors.
    \item Any non-trivial contribution by a generative-AI tool must be disclosed in the relevant notebook or script, as specified in the course manual.
    \item The rules on generative-AI use in the course manual apply in full. In particular, students may not upload course datasets or full analysis notebooks to external generative-AI services, or use generative AI as a black-box analysis engine for graded work.
    \item Students are free to select any topic of interest---including topics outside traditional communication science or the social sciences---provided that it is framed around a research question and based on original data analysis.
\end{itemize}

\subsection*{Submission instructions}

Submit both the PDF report and the \texttt{.zip} folder containing code and outputs via Canvas.

Deadline: Submission is open until 17 December, 23:59. Late submissions are handled in accordance with the course manual and applicable examination regulations.
